% Source: https://artofproblemsolving.com/community/c3h2297782_easy_mathz
% Stable identity: aops-topic-2297782-q01
A student finds the average of three scores $x$, $y$, $z$ in two ways: as $\dfrac{x+y+z}{3}$, and by first averaging $x$ and $y$ and then averaging the result with $z$, that is, $\dfrac{\dfrac{x+y}{2}+z}{2}$. For which values of $x$, $y$, $z$ do the two methods give the same result?
